\section{Introducción}
En la era digital, la gobernanza y la auditoría de sistemas se han convertido en pilares fundamentales para garantizar la continuidad, seguridad y alineación estratégica de la tecnología con los objetivos organizacionales. Las instituciones de salud, en particular, enfrentan el reto de digitalizar sus operaciones críticas de manera acelerada, a menudo introduciendo vulnerabilidades complejas en sus entornos locales y remotos.

Este documento presenta una propuesta técnica de auditoría integral orientada a evaluar y robustecer la infraestructura, controles y gobernanza de la Clínica Médica Andaluz S.R.L., integrando de manera sinérgica los marcos internacionales COBIT 2019, ISO/IEC 27001 e ITIL 4.

\section{Objetivos}

\subsection{Objetivo General}
Diseñar una propuesta profesional de Auditoría Integral para la Clínica Médica Andaluz S.R.L., que permita evaluar la efectividad de sus controles de red, seguridad lógica en accesos remotos y gobernanza de TI, garantizando la co-creación de valor y la mitigación de riesgos de seguridad.

\subsection{Objetivos Específicos}
\begin{itemize}
    \item Diagnosticar el estado actual de los accesos remotos al sistema de Historia Clínica Electrónica (HCE).
    \item Evaluar el diseño y segmentación de la infraestructura de red física y lógica local.
    \item Proponer controles basados en ISO/IEC 27001 para mitigar las vulnerabilidades identificadas.
    \item Definir un esquema de gobernanza basado en COBIT 2019 y un modelo de entrega de servicios alineado con ITIL 4.
    \item Planificar la ejecución de la auditoría y sus respectivas actividades mediante un marco ágil de gestión.
\end{itemize}
